% Chapter 9b: Revised Protocols
% Scene function: Mandatory training. The system absorbing time.
% Kafka comedy. Meri asks a question Compliance cannot answer.

\chapter{Revised Protocols}

On Wednesday there is a mandatory training. The subject line on the calendar invitation says ``Revised Protocols for Environmental Documentation in Non-Standard Dwelling Types.'' Attendance is required. The invitation was sent by Compliance.

The training is in the large conference room, which is the room used for events that are too important to skip and too unimportant to prepare for. The chairs are arranged in rows. There is a projector. The projector displays a title slide that contains the full name of the training, the date, a Compliance logo, and the word ``Draft'' in the lower right corner. The title slide has been on screen for ten minutes. Nobody has advanced past it.

The Compliance officer arrives. They apologize for the projector, which is working correctly. They advance to the second slide. It is a flowchart.

The flowchart describes the process by which a field operative should document environmental conditions in a dwelling that does not conform to standard residential classification. It has fourteen boxes and twenty-three arrows. Several of the arrows point to boxes that say ``See Supplemental Flowchart B.'' Supplemental Flowchart B is not in the slide deck. The Compliance officer mentions this and says it will be distributed separately.

I look around the room. There are eleven operatives present. Two are from Possessions. One is from Whispers. The rest are Hauntings. Nobody is taking notes except Meri, who is in the second row with her pen out, writing in what appears to be a new notebook. She has graduated from the margins of the Operations Manual.

The Compliance officer spends fifteen minutes explaining a change to Field 14 on Form 7-C. Field 14 previously required the operative to record the dwelling's construction material. It now requires the operative to record the dwelling's construction material and its acoustic classification, which is a category I have never seen on a form before and which, if it had existed six years ago, would have saved me the trouble of writing ``good acoustics'' on a clipboard during a Week 1 Field Assessment.

I have been filling in Field 14 without the acoustic classification for six years. Either I have been making an error nobody noticed, or the field has been redefined and my previous entries are now retroactively incomplete. The Compliance officer does not address which of these is the case.

A Possessions operative raises a hand and asks what acoustic classification means in the context of a human body, which is where Possessions does most of its work. The Compliance officer says that is an excellent question and refers them to the supplemental flowchart.

Meri raises her hand.

``If the acoustic classification affects the Fear Index calculation, should we retroactively adjust the baseline projections for active cases?''

The room is quiet for a moment. The Compliance officer looks at the flowchart. The flowchart does not contain an answer to this question. The flowchart does not contain the word ``retroactively.''

``That would be a question for your supervisor,'' the Compliance officer says.

The training ends. I have been in this room for forty-five minutes. I know one new thing: Field 14 on Form 7-C now has a second component. I will fill it in. I have always filled in the fields.
